\subsection{Semi-simplicity of finite-dimensional representations}

We recall (INT, VII, p. 71, § 3, n° 1, lemma 1) that for every continuous representation \(\varrho\) of G in a Hilbert space E, there exists a nondegenerate positive Hermitian form \(q\) on E such that the topological vector space structure of E defined by q is identical to the initial structure of E, and such that \(\varrho\) is a unitary representation of G in the Hilbert space E equipped with the inner product q.

PROPOSITION 1. --- Let \(\varrho\) be a finite-dimensional linear representation of G in a separated topological vector space E. There exists an inner product \(q\) on E such that \(\varrho\) is a unitary representation in the Hilbert space E equipped with \(q\) . In particular, \(\varrho\) is semisimple.

Since E is finite-dimensional, there exists a Hilbert space structure on E. The result follows by applying the preceding remark and Corollary 2 of V, p. 392.

Remark. --- We shall see later (Cor. 4 of V, p. 466) that every unitary representation of G is semi-simple.

COROLLARY. --- Let \(\varrho_{1}\) and \(\varrho_{2}\) be finite-dimensional representations of G. The representations \(\varrho_{1}\) and \(\varrho_{2}\) are isomorphic if and only if their characters are equal.

This follows from the proposition and Corollary 3 of V, p. 392.

